Chương này trình bày các cơ sở lý thuyết liên quan sử dụng trong đồ án. Nội dung 2.1, 2.2 tổng hợp các kiến thức cơ bản về cơ học vật rắn, bao gồm mô hình ma sát, động học–động lực học vật thể, bài toán quy hoạch toàn phương và các khái niệm cơ bản về grasping. Phần 2.3, 2.4 và 2.5 trình bày các thuật toán tối ưu cơ bản dùng làm tiền đề cho phương pháp nghiên cứu được tham khảo từ nhóm nghiên của thuộc đại học Stuttgart \cite{Rosenfelder2024}   

\section{Mô hình động học-động lực học của vật thể rắn}
Vật thể cần được thao tác được mô hình như một vật rắn tuyệt đối (rigid body) tức là có hình dạng cố định, không biến dạng dưới tác động của lực, có vị trí $\mathbf{x}(t) = [x(t)\; y(t)]^{\top} \in \mathbb{R}^2$ trong đó $x(t)$ và $y(t)$ lần lượt là toạ độ trên hai trục x và y, góc quay $\boldsymbol{\theta}(t)$. Chuyển động của vật được coi như là sự tổng hợp chuyển động tịnh tiến được mô tả bởi gia tốc tịnh tiến $\mathbf{\ddot{x}}$ và chuyển động quay được mô tả bởi gia tốc góc $\boldsymbol{\ddot{\theta}}$ tại điểm trọng tâm của vật trên mặt phẳng hai chiều. 

Mô hình ma sát Coulomb giả định rằng lực ma sát chỉ phụ thuộc vào lực pháp tuyến $\mathbf{F}_{n}$ và hệ số ma sát giữa hai bề mặt, và không phụ thuộc vào diện tích tiếp xúc hay vận tốc chuyển động, ngoại trừ chiều của vận tốc tịnh tiến $\dot{\mathbf{x}}$ và vận tốc góc $\boldsymbol{\dot{\theta}}$. Khi vật chưa trượt, lực và mô men ma sát tĩnh điều chỉnh trong giới hạn $\|\mathbf{F}_{f}\|\le \mu_{t} \mathbf{F}_{n}$ và $\left| \mathbf{M}_{f} \right| \le \mu_{r} \mathbf{F}_{n}$ để chống lại chuyển động tương đối, với $\mu_{t}$ và $\mu_{r}$ lần lượt là hệ số ma sát tĩnh tịnh tiến và hệ số ma sát tĩnh quay. Khi trượt xảy ra, lực ma sát động có độ lớn không đổi và ngược hướng vận tốc tiếp tuyến, được mô tả bởi:
\begin{align}
    \mathbf{F}_{f} = -\mu_{t}' \mathbf{F}_{n} \frac{\dot{\mathbf{x}}}{\|\dot{\mathbf{x}}\|}
    \\[6pt]
    \mathbf{M}_{f} = -\mu_{r}' \mathbf{F}_{n}\,\operatorname{sgn}(\boldsymbol{\dot{\theta}})
\end{align}
trong đó $\mu_t'$ và $\mu_r'$ lần lượt là là hệ số ma sát động tuyến tính mà xoay của vật, hàm $\operatorname{sgn}$ là một hàm dấu.

Khi vật chịu tác động của $n$ lực đẩy ngoài $\mathbf{f}_i \in \mathbb{R}^2$ với các cánh tay đòn tương ứng $\mathbf{r}_i \in \mathbb{R}^2$, ta có thể mô tả động lực học của hệ bằng cách áp dụng định luật~II Newton cho chuyển động tịnh tiến và phương trình Euler cho chuyển động quay. Cụ thể, hệ phương trình động lực học được viết dưới dạng:
\begin{align}
    m \ddot{\mathbf{x}} &= \sum_{i=1}^{n} \mathbf{f}_i + \mathbf{F}_{f} ,
    \\[6pt]
    I \boldsymbol{\ddot{\theta}} &= \sum_{i=1}^{n} ( \mathbf{r}_i \times \mathbf{f}_i) + \mathbf{M}_{f}
    \label{eq:ang}
\end{align}
trong đó $m$ và $I$ lần lượt là khối lượng và mô men quán tính của vật.